% ============================================================================
%  DRAFT rewrite of Section II (absorbing Section III), sphere-kernel model.
%  2026-07-20.  Figure: fig3_model.pdf (/tmp/fig3_model.png).  Not yet merged
%  into main.tex.  Equation/figure \labels are placeholders.
% ============================================================================
\section{The amplification model}\label{sec:model}

The model is a single algebraic source term added to the Spalart--Allmaras
transport equation, built entirely from the local mean-velocity profile. This
section constructs it from what it must reproduce, on a geometric picture of
the local profile that we build up in turn.

\subsection{What the kernel must reproduce}\label{sec:reproduce}
The canonical laminar instabilities are each a self-similar analytical
profile: the attached Falkner--Skan family, from favorable gradients through
the Blasius layer (viscous Tollmien--Schlichting) to the incipient-separation
limit $H\!=\!4.03$; the reversed-flow lower branch (Stewartson), i.e.\ the
separation-bubble shear layer; and the free shear layers---mixing layer, jet,
and wake. The constant-curvature, zero-wall-shear \emph{parabola}
(stagnation profile) is the neutral reference. For every one of these, linear
stability tells the kernel \emph{two} things it must get right: not only the
growth \emph{rate}, but \emph{where}, across the layer, the disturbance
extracts its energy---the critical layer / production core. The second is
easy to overlook and, as we now argue, decisive.

\subsection{Temporal versus spatial growth: why location matters}\label{sec:tempspatial}
The present source is a \emph{temporal} rate: it multiplies the working
variable, $\partial_t\tilde\nu = a\,\Omega\,\tilde\nu$ (per unit time), and is
carried downstream by the local mean velocity,
$u\,\partial_x\tilde\nu = a\,\Omega\,\tilde\nu$. Transition, however, is set by
the \emph{spatial} amplification---the $N$-factor accumulated along the
wall. The two differ by the local advection speed: the spatial growth the
model deposits at height $y$ is $(a\,\Omega)/u(y)$. Placing the rate at the
wrong wall-normal station therefore does not merely mis-scale the growth---%
through the $1/u$ weighting it can \emph{over}-amplify wherever $u$ is small.
The laminar separation bubble is the sharp case: the physical wavepacket rides
the detached shear layer at $u\!\approx\!c\!>\!0$, but a kernel that lets
growth slide toward the dividing streamline ($u\!\to\!0$) deposits a diverging
spatial rate---the dead-air over-amplification that drives $N$ past $e^{9}$.

This corrects a premise of the earlier drafts, which calibrated the rate as if
only its magnitude mattered. The kernel must \emph{resolve where the
instability grows}, not only how fast. This is what motivates a local geometry
rich enough to separate ``which state'' the profile is in from ``how much'' it
should amplify.

\subsection{The indicator sphere}\label{sec:sphere}
To second order about any wall-normal point, the profile is fixed by three
velocity scales formed from the wall distance $d$: the velocity $u$, the shear
$d\,u'$, and the curvature $\tfrac12 d^2 u''$ (the Taylor coefficients weighted
by $d^{k}$; the $\tfrac12$ is the quadratic term's actual contribution to the
velocity). Dividing by their common magnitude $R$ places the local state on
the unit sphere of directions
\begin{equation}
  (X,Y,Z) = \big(u,\; d\,u',\; \tfrac12 d^2 u''\big)\big/R,
  \qquad R=\sqrt{u^2+(d u')^2+(\tfrac12 d^2u'')^2}.
  \label{eq:indicators}
\end{equation}
Because the growth is invariant under $u\!\to\!-u$ (a layer and its mirror are
equally unstable), the state is a direction defined only up to sign---a point
of the real projective plane $\mathbb{RP}^2$, the space of lines through the
origin in $\mathbb{R}^3$.\footnote{$\mathbb{RP}^2=S^2/\{\pm\}$; the sphere
double-covers it. We plot one hemisphere, with the antipode of any point
understood to be identified with it.} We call $X,Y,Z$ the velocity, shear, and
curvature indicators; the shear/velocity balance is the profile-fullness
$\Gamma = 2Y^2/(X^2+Y^2)\in[0,2]$ used in earlier work.

\subsection{Where the instabilities live}\label{sec:live}
Tracing each analytical profile of \S\ref{sec:reproduce} as a curve on
$\mathbb{RP}^2$, and thickening the wall-normal band where its production
peaks, gives Fig.~\ref{fig:model}. The result is uniform: \emph{every}
amplifying band lies on the shear-dominated side, toward the shear pole $+Y$
($\Gamma\!\to\!2$), and away from the velocity pole (freestream) and from
pure-curvature states. The stagnation profile sits exactly on the dividing
curve. The region the kernel must switch on is thus bounded below by the
parabola and, on its flanks, by the profile-fullness $\Gamma$---an empirical
statement we now make exact.

\begin{figure}[tp]
  \centering
  \includegraphics[width=0.9\textwidth]{fig3_model.pdf}
  \caption{What the kernel must reproduce. The analytical profiles on the
  $\mathbb{RP}^2$ indicator sphere \eqref{eq:indicators} (x-ray view down the
  velocity$=$shear diagonal; curvature vertical, shear$-$velocity horizontal;
  solid$=$near hemisphere, dotted$=$far). Each profile is a thin curve; the
  \textbf{thick} band marks where its linear-stability production peaks---where
  the disturbance actually grows. All bands lie on the shear side of the
  parabola great circle $g=0$ (bold), toward the shear pole. The stagnation
  (pure-parabola) profile lies on $g=0$ and is neutral.}
  \label{fig:model}
\end{figure}

\subsection{The parabola great circle and the Rayleigh coordinate}\label{sec:parabola}
The wall-bounded parabolas $u = By + Cy^2$ form \emph{exactly} a great circle:
for every $B,C$ the trajectory \eqref{eq:indicators} satisfies the single
linear relation
\begin{equation}
  g \;\equiv\; Y - X - Z \;=\; \text{shear} - \text{velocity} - \text{curvature} \;=\; 0 .
  \label{eq:parabola}
\end{equation}
Off the parabola, $g$ has a direct meaning:
$g(d) = -\int_0^d \tfrac12 s^2\,u'''(s)\,\mathrm{d}s$, the accumulated
curvature reversal, a local measure of \emph{inflectional content}. A parabola
has no inflection ($u'''\!=\!0\Rightarrow g\!=\!0$) and by Rayleigh's
inflection theorem is inviscidly neutral; $g>0$ is the inflectional, unstable
side. Thus $g$ is the sphere's Rayleigh coordinate and the parabola great
circle \eqref{eq:parabola} is the neutral boundary---the lower edge of the
amplifying region seen in Fig.~\ref{fig:model}.

\subsection{The rate}\label{sec:rate}
Two sphere functions build the rate. The production weight is the shear
fraction
\begin{equation}
  \hat S \;=\; \frac{Y}{\sqrt{X^2+Y^2}} \;=\; \sqrt{\Gamma/2},
  \label{eq:shat}
\end{equation}
which encodes that production scales with $|U'|$. Crucially $\hat S$ is a
\emph{ratio}---it depends on longitude only---so it stays finite when the
curvature indicator, carrying its $d^2$ scale, dominates far from the wall: a
detached shear layer driven onto the curvature pole still registers through its
$\Gamma$, and is distinguished there from a no-shear velocity extremum
($\Gamma\!\to\!0$). The rate is
\begin{equation}
  a \;=\; \operatorname{clip}\big\langle \hat S\,G(g)\big\rangle_{+},
  \qquad G(g) = g + c_{\!\sqrt{}}\sqrt{g},\quad c_{\!\sqrt{}}\!\approx\!0.13,
  \label{eq:rate}
\end{equation}
with the ceiling $a_{\max}\!\approx\!0.21$ (the free-shear-layer eigenvalue,
\S\ref{sec:calib}). The monotone shaping $G(g)$ is coined by calibration
(Fig.~\ref{fig:calibrate}a): the physical growth-per-vorticity follows
$\propto g$ at large $g$ but $\propto\!\sqrt{g}$ near the parabola---the
sublinear near-neutral behavior confirmed independently by the Drela--Giles
envelope, $\mathrm{d}N/\mathrm{d}Re_\theta\!\propto\!g^{0.66}$ with a finite
Blasius rate. A purely linear $G(g)=g$ would under-amplify Blasius by a factor
of two.

\subsection{The onset}\label{sec:onset}
The rate \eqref{eq:rate} is scale-free; the viscous inception is a separate,
multiplicative gate rising smoothly from $0$ to $1$ across the neutral point,
\begin{equation}
  S\!\big(Re_\Omega / Re_\Omega^{\mathrm c}(g)\big),
  \qquad Re_\Omega = d^2\Omega/\nu .
  \label{eq:onset}
\end{equation}
Its threshold is a function of the sphere location: reading $Re_\Omega$ at each
profile's neutral $Re_\theta$ (Fig.~\ref{fig:calibrate}b) collapses onto a
monotone $Re_\Omega^{\mathrm c}(g)$, falling from $\approx\!450$ at the
parabola (Blasius, the high viscous critical Reynolds number,
$Re_\Omega^{\mathrm c}\!=\!2.19\,Re_\theta^{\mathrm c}$) to $\approx\!60$ for
inflectional layers (Kelvin--Helmholtz, essentially immediate). Rate and onset
reinforce: near the parabola, low rate \emph{and} high threshold give the late,
slow transition of a zero-pressure-gradient layer; inflectional profiles get
high rate \emph{and} low threshold, hence early, fast transition. The complete
source is
\begin{equation}
  P_{\mathrm{AI}} \;=\; a_{\max}\,\hat S\,G(g)\;
  S\!\big(Re_\Omega/Re_\Omega^{\mathrm c}(g)\big)\;\Omega\,\tilde\nu .
  \label{eq:pai}
\end{equation}
Two sphere functions---the rate $\hat S\,G(g)$ and the onset threshold
$Re_\Omega^{\mathrm c}(g)$---replace the four-factor construction of earlier
drafts (gate $Q$, switch $S(z)$, favorable-rate factor $f_\lambda$, ceiling).
Favorable-gradient suppression is no longer a separate factor: a favorable
gradient carries the state to $g<0$, where the rate is off by construction.

\subsection{Calibration}\label{sec:calib}
[Absorbs former \S III. $a_{\max}$ from the inflectional family; $G(g)$ and the
$\sqrt g$ coefficient from the growth-per-vorticity fit; $Re_\Omega^{\mathrm
c}(g)$ from the neutral-point collapse; freestream-turbulence scaling of the
onset. Figure~\ref{fig:calibrate}: (a) rate, log-log with the $\sqrt g$/$g$
references and the Drela--Giles envelope; (b) onset collapse.]
